\cotefile{16}
\scanpage{17}\mspage{14}
\uline{Cor} Soit $E$ un anneau gradué, \inserted{$2n$ un entier pair $\geq 0$,} et soit $L\in E^{(2)}$. S'il existe un hom d'anneaux gradués
\[
u\colon \unsure{\Phi|\mathbb{\Phi}}\longrightarrow E
\]
tel que $u(\mathbb L)=L$, et tel que les $u(\Pi_i)=\pi_i$ définissent la graduation de $E$ par (1.3), alors un tel $u$ est uniquement déterminé, et les $\pi_i$ sont déterminés de façon unique, en termes de la structure \inserted{d'anneau} gradué de $E$, par les conditions (1.1) et (1.3). \deleted{\unsure{En termes des} $\pi_i$,} \inserted{Si on a \inserted{déjà} des $\pi_i$ satisfaisant ces conditions,} l'existence de $u$ équivaut alors à \deleted{l'\unsure{existen|existence}} la validité, pour \deleted{\unsure{tt}} $0\leq i\leq n$, de la condition
\[
(\quad\exists\,w_i\in E^{-(2n-2i)}\quad\text{tel que}\quad (w_i L^{n-i}-1)\pi_{2n-i}\deleted{\unsure{=0}}\unsure{\}}\;=0,
\]
et \deleted{$u(\unsure{\Lambda})$} le système des $\lambda=u(\Lambda)$, \deleted{et} $\pi_{i,\alpha}=u(\Pi_{i,\alpha})$ est alors caractérisé (en termes de $L$ et des $\pi_i$) par les équations (6.14). \deleted{\unsure{La condi}} \inserted{Ou encore\unsure{:} $\exists$ un module fidèle $M^*$ sur $E^*$ tel que, pour la graduation conjuguée \unsure{i.e.} ($2n-p$ au lieu \unsure{de $p$})}

\unsure{Indépendamment} de la donnée des $\pi_i$, l'existence de $u$ équivaut à l'existence d'un module \uline{gradué} fidèle sous $E^*$, \deleted{tel que} \deleted{\unsure{les}} (d'où les $\pi_i\in\operatorname{End}(M^*)$), et tel que \ldots\ldots .

\uline{Question} Si on ne se donne \uline{pas} la graduation de $E$ à priori, \inserted{(elle sera définie à posteriori par les $\pi_i=u(\Pi_i)$)} reste-t-il vrai qu'il y a au plus un hom $\Phi\to E$ tel que $u(\mathbb L)=L$ ?

\textcircled{Non}, déjà dans le cas $n=1$, $2n=2$, où il y a une base de $\Phi$ formée de $\pi_0,\pi_1,\pi_2,L,\Lambda$, et où $\mathbb Z[L]$ n'est que $\mathbb Z+\mathbb Z L$ ! Il faut donc se donner en même temps \uline{$L$ et les $\pi_i$}.
% archive sheet number: 16
\scanpage{18}
8) Soit maintenant $E^*$ un anneau gradué, \deleted{\unsure{muni de}} tel qu'il existe $L\in E^{(2)}$ satisfaisant aux \uline{conditions de Lefschetz} (\deleted{de} les conditions du corollaire du n° précédent). On dit que $E^*$ est un \textquotedblleft anneau \deleted{d'\unsure{opérateurs} lefschetziens \unsure{(à ne pas confondre avec la notion}} \inserted{de correspondances lefschetz\unsure{ien|iens}\textquotedblright. (Attention,} ce n'est \unsure{pourtant} qu'un \unsure{résidu} de Lefschetz), et $L$ en est un \textquotedblleft élément de Lefschetz\textquotedblright. \deleted{(\unsure{on dira simplement})\unsure{,}} L'élément $\Lambda$ associé est \deleted{appelé} \textquotedblleft l'élément \deleted{\unsure{adjoint formel}} \deleted{\unsure{anti}} \inserted{co.}lefschetz\deleted{\unsure{ien}} \textquotedblright\ associé à $L$. (NB. À montrer que $L$ se déduit de $\Lambda$ comme $\Lambda$ de $L$, par simple \unsure{symétrie|symétrique} sur les degrés de $E^*$ \deleted{$\unsure{\gtrless}-$}, et passage à l'anneau opposé --- sorte de \deleted{\illegible{1}} transposition).
% slanted
\margin{ou mieux : \textquotedblleft principe de dualité\textquotedblright\ pour les anneaux de correspondances lefschetziens.}

On peut se demander de \uline{caractériser} les \deleted{\unsure{les}} éléments de Lefschetz $L$. Une condition nécessaire, lorsque \unsure{la condition de|la condition sur} \unsure{$2n$} n'est pas \unsure{rédondant|rédondante}, i.e. que
% slanted
\margin{\unsure{désignons par|déjà par prop.|déjà pour $n$}}
\[
\begin{aligned}
\overset{\unsure{1)}}{E^{(2n)}}\ne0\quad&\deleted{\unsure{\Leftrightarrow}}\Leftrightarrow\overset{1')}{E^{(-2n)}}\ne0\deleted{\unsure{\}}}
\quad\Leftrightarrow\overset{2)}{\pi_0},\deleted{\unsure{\pi_{2n}}}\ne0\deleted{\unsure{\}}}\\
&\inserted{\Leftrightarrow\overset{2')}{\pi_{2n}}\ne0}
\Leftrightarrow\overset{3)}{L^n}\ne0\quad\Leftrightarrow\overset{3')}{\Lambda^n}\ne0\quad\deleted{\Leftrightarrow}.
\end{aligned}
\]
c'est que l'on ait $L^n\ne0$. On a parfois une réciproque, dans le cas où ($E$ étant une \uline{algèbre} sur un anneau commutatif $k$)
% archive sheet number: 17
\scanpage{19}\mspage{45}
\deleted{en groupes \deleted{localement}\inserted{lisse et} de type fini sur $k$, soit $\underline{\mathrm{Pic}}_U$, et un isomorphisme
\[
\mathrm{Pic}(U)=\underline{\mathrm{Pic}}_U(k) \tag*{(8.13)}
\]
de $\mathrm{Pic}(U)$ avec le groupe des points à valeurs dans $k$ de $\underline{\mathrm{Pic}}_U$. Cela nous fournit donc pour $\mathrm{Pic}(U)$ une structure d'extension
\[
0\handwritten{\to}\mathrm{Pic}(U)^{\unsure{o|0}}\handwritten{\longrightarrow}\mathrm{Pic}(U)\handwritten{\longrightarrow}\mathrm{NS}(U)\handwritten{\longrightarrow}0\ , \tag*{(8.14)}
\]
où on a posé
\[
\mathrm{Pic}(U)^o=\underline{\mathrm{Pic}}_U^o(k)\ ,\quad\mathrm{NS}(U)=\underline{\mathrm{NS}}_U(k)\ , \tag*{(8.15)}
\]
\deleted{\unsure{où}} $\underline{\mathrm{Pic}}_U^o$ désignant la composante neutre de $\underline{\mathrm{Pic}}_U$, et $\underline{\mathrm{NS}}_U$ le quotient de $\underline{\mathrm{\deleted{\unsure{B}}Pic}}_U$ par cette composante. La construction indiquée dans loc. cit. permet d'ailleurs de préciser la structure de $\mathrm{NS}(U)$, en introduisant la matrice d'intersection
\[
(\mathrm C_i.\mathrm C_j)_{1\leq i,j\leq r}
\]
définie par les $r$ composantes irréductibles $\mathrm C_i$ de $f^{-1}(x)_{\text{réd}}$ sur le \margin{\handwritten{$\times$}}schéma régulier $X'$ : en vertu de MUMFORD\inserted{\handwritten{[28]}}, cette matrice est définie néggtive, donc le conoyau de l'homomorphisme
\[
\underline{\mathrm Z}^{r}\handwritten{\longrightarrow}\underline{\mathrm Z}^{r}
\]
qu'elle définit est un groupe fini, dont l'o\deleted{\unsure{dr}}rdre n'est autre que la valeur abso\deleted{\unsure{h}}lue du déterminant de cette matrice. C\deleted{\unsure{o}}eci posé, on trouve que $\mathrm{NS}(U)$ est canoniquement isomorphe à ce groupe fini. La matrice envisagée étant symétrique, on voit aussitôt que ce groupe est canoniquement autodual (par un accouplement \deleted{\unsure{parfait}}\handwritten{symétrique} à valeurs dans $\underline{\mathrm Q}/\underline{\mathrm Z}$), d'où une autodualité\inserted{\handwritten{ symétrique}}
\[
\mathrm{NS}(U)\mathbin{\mathrm x}\mathrm{NS}(U)\handwritten{\longrightarrow}\underline{\mathrm Q}/\underline{\mathrm Z}\ , \tag*{(8.16)}
\]
dont on constate aisément qu\deleted{e}'elle est indépendante de la résolution $X'$ choisie.

Si maintenant \handwritten{$\unsure{l|\ell}$} est un nombre premier distinct de $\mathrm{car}\ k$, alors la structure connue des groupes algébriques commutatifs connexes sur $k$ \margin{\handwritten{$\times$}}\inserted{\handwritten{[9, exp.4]}} mont\deleted{\unsure{e}}re que $\mathrm{Pic}(U)\mathrm{o}$ est un groupe \handwritten{$l$}-divisible, donc}
% archive sheet number: 18
\scanpage{20}
\[
\begin{aligned}
\deleted{\unsure{b'}}\quad b_0\quad&
\underset{\rule{.35pt}{.55ex}\rule{4.2em}{.35pt}\rule{.35pt}{.55ex}}{b_1\quad b_2\quad b_3}\quad b_4\\
n b_0+(n-1)b_1={}&\hspace{6em} L\\
\underset{\rule{.35pt}{.55ex}\rule{4em}{.35pt}\rule{.35pt}{.55ex}}{2b_0+b_1}&\quad 2
\end{aligned}
\]
\[
\lambda^{(n-i)b_i}\qquad \mathcal E^{n-i}\qquad
\unsure{\mathscr E|E}^{\unsure{i|\bullet}}\times\mathscr E^{\unsure{2n-2\cdot|2n-2i}}
\]
$\deleted{\unsure{\lambda}}$
\diagram{\begin{tikzcd}[ampersand replacement=\&] M^0\arrow[r]\&M^{2n}\end{tikzcd}}
\[
\begin{aligned}
&\unsure{*}\quad\deleted{\Lambda\quad\lambda L}\unsure{\leadsto}\\
& L\unsure{\leadsto|\sim}\lambda L\quad\text{\unsure{alors|Alors}}\quad \Lambda\leadsto\tfrac1\lambda\Lambda
\end{aligned}
\]
\[
\begin{aligned}
\Lambda&=P(L\quad\deleted{\unsure{\xi}} &&\deleted{\unsure{LL'}\quad {L'}^{2}=LL'}\\
\Lambda&=\Omega(L)\,\Delta(L)^{-N}. &&\underline{L'}\quad\underline{LL'}\quad\underline{L^{\unsure{2|3}}L'}\quad\deleted{L^4}\quad L^3L'\\
&&&\bigl(L^4,\ {L'}^{2}-LL',\ L^2L'-L^3,\\
&&&\hspace{11em} L^3L'\bigr)\\
&&&1\ \big|\ L,L'\ \big|\ L^{\unsure{2|?}},LL'\ \big|\ L^3\ \bigr]
\end{aligned}
\]
\diagram{\begin{tikzcd}[ampersand replacement=\&,column sep=small]
H^0\arrow[rrrr,bend right=20]\& H^1\arrow[rr]\& H^2\&H^3\& H^4
\end{tikzcd}}
\[
\begin{aligned} {L'}^{3}\deleted{\unsure{\xi}}&={L'}^{2}L=L^2L'\\&=L^3\end{aligned}
\]
\diagram{\begin{tikzcd}[ampersand replacement=\&,row sep=small]
H^0\arrow[r]\&H^2\arrow[r,dash]\&H^4\arrow[r]\&H^6\\
\text{\uuline{rg 1}}\&\text{\uuline{rg 2}}\&\text{\uuline{rg 2}}\&\text{\uuline{rg 1}}
\end{tikzcd}}
\[
\begin{aligned}
1\quad\deleted{\unsure{\Sigma V}}\quad &\operatorname{Sym}^{\unsure{\bullet|3}}(V)\; \unsure{-|{\to}}\;V-\\
&\unsure{k}\\
L,L'\qquad &L^2=\overline L,\ L\overline{L'}=\overline{L'},\ \overline{L'}^{\,2}=\overline{L'}\\
& L^{\unsure{2|4}},L',LL'.\qquad L'\\
L,L'&\\
L^2,LL'\quad & {L'}^{2}=LL'\qquad {L'}^{3}=L'{L'}^{2}=L^2L'\\
L^3=L^2L'\quad &\mathbb Z[L]/(L^4)\deleted{\unsure{E}}\,[L']/{L'}^{2}-L^{\prime\prime}
\end{aligned}
\]
% archive sheet number: 19
\scanpage{21}\mspage{42}
on trouve une suit\deleted{\unsure{é}}e exacte canonique \handwritten{(où $\mathbb Z_l[1]=\varprojlim_{\nu}\deleted{\unsure{Z}}\mu_{l^\nu}$ est le faisceau $l$-adique de TATE) :}
\[
0\handwritten{\longrightarrow}t^{i-1}(X,l)\handwritten{\longrightarrow}H^i(X,\underline{\mathrm Z}_l\handwritten{[1]})\handwritten{\longrightarrow}T_l(H^i(X,\handwritten{\mu}_{l^\infty}))\handwritten{\longrightarrow}0\ , \tag*{(8.3)}
\]
où \deleted{$\unsure{t|H}^{i-1}(X,l)$} on a posé
\[
t^{i-1}(X,l)=H^{i-1}(X,\handwritten{\mu}_{\deleted{\unsure{l}}l^\infty})/H^{i-1}(X,\handwritten{\mu}_{l^\infty})^{\unsure{o|0}} \tag*{(8.4)}
\]
qui est donc un groupe fini de $l$-torsion, et s'identifie donc, par (8.3), au sous-\deleted{groupe}\inserted{module} de torsion d\deleted{\unsure{a|u}}\unsure{u|a} $\underline{\mathrm Z}_l$-module de type fini
\[
H^i(X,\underline{\mathrm Z}_l\handwritten{[1]})=\lim H^i(X,\deleted{\unsure{\underline{\mathrm Z}/l^\nu}}\handwritten{\mu_{l^\nu}})\ , \tag*{(8.5)}
\]
dont la partie libre s'identifie, à son tour, à $T_l(H^i(X,\handwritten{\mu}_{l^\infty}))$. On en \margin{\handwritten{!}}déduit également un isomorphisme
\[
\deleted{\unsure{H}}(H^i(X,\handwritten{\mu}_{l^\infty})^o\handwritten{\simeq} H^i(X,\underline{\mathrm Z}_l\handwritten{[1]})\otimes_{\underline{\mathrm Z}_l}\underline{\mathrm Q}_l\ /\ \mathrm{Im}\ H^i(X,\underline{\mathrm Z}_l\handwritten{[1]}), \tag*{(8.6)}
\]
qui montre que $H^i(X,\handwritten{\mu}_{l^\infty})$ et $H^i(X,\underline{\mathrm Z}_l\handwritten{[1]})$ se déterminent mutuellement mod groupes \deleted{\unsure{dont}} finis.

\handwritten{8.2.} Un passage à la limite analogue dans l\deleted{es}a suite\deleted{s} exacte\deleted{s} de Kummer (BR II 3.1) nous donne de même une suite exacte
\[
0\handwritten{\longrightarrow}\mathrm{NS}(X)\otimes_{\underline{\mathrm Z}}\underline{\mathrm Z}_l\handwritten{\longrightarrow}H^2(X,\underline{\mathrm Z}_l[1])\handwritten{\longrightarrow}T_l(\mathrm{Br}'(X))\handwritten{\longrightarrow}0\ , \tag*{(8.7)}
\]
en supposant maintenant que $\mathrm{Pic}(X)$ est extension d'un groupe de type fini $\mathrm{NS}(X)$ (le \uline{groupe de Néron-Sévéri}) par un groupe divisible, et que $H^2(X,\underline{\mathrm Z}/l\underline{\mathrm Z})$ est fini. \deleted{\unsure{(}}On trouve en passant le fait, classique dans la théorie transcendante, que \deleted{le}\inserted{la composante $l$-primaire du} sous-groupe de torsion de $\mathrm{NS}(X)$ est is\deleted{\unsure{o}}omorphe au sous-groupe de t\deleted{\unsure{ro}}orsion de $H^2(X,\underline{\mathrm Z}_l[1])$, en même temps qu'on précise du point de vue de la cohomologie $l$-adique la relation entre le groupe de Brauer (ou plus précisément, $T_l(\mathrm{Br}'(X))$) et la\textquotedblleft partie transcendante\textquotedblright\ du $H^2$ \handwritten{,} \deleted{\unsure{;}}comparer (BR II 3 passim ).

Lorsque $X$ est \inserted{une surface} propre, lisse et connexe sur un corps $k$ algébriquement clos, il y a lieu de considérer sur le terme médian de (8.7) la forme bilinéaire symétrique canonique, fournie par le cup-produit,\deleted{\unsure{à}} qui est à valeurs d\deleted{\unsure{q}}ans
% archive sheet number: 20
\scanpage{22}
$E^{2n}$ \deleted{est engendré par}\inserted{admet la \uline{base}} $L_0^n$, \deleted{$E^{-2n}$ engendré par $\Lambda_0^n$, pour \uline{un} choix de $(L,\Lambda)$. Alors pour tt \unsure{$x$}$\in E^{2n}$, on a \unsure{$L=L'$}} \inserted{(pour un $L_0$ de Lefschetz fixé)} \unsure{Or} --- \unsure{où|on}, si $L'$ est un autre élément Lefschetz(ien)
\[
\begin{aligned}
L_0^n&=L_0^n\pi_0=\underset{\rule{.35pt}{.5ex}\rule{3.1em}{.35pt}\rule{.35pt}{.5ex}}{L_0^n{\Lambda'}^{n}}{L'}^{n}
=L_0^n(\beta\Lambda_0^n)(\deleted{\unsure{\beta}}\alpha L_0^n)\\
&=L_0^{\unsure{+n|n}}\alpha\beta\,\pi_0=\alpha\beta\, L_0^n.
\end{aligned}
\]
ce qui prouve que $\alpha\beta=1$, donc $\alpha$ et $\beta$ inversibles, donc \deleted{\unsure{il est}} $L_0^{\unsure{+n|n|\prime n}}$ et $\Lambda_0^{\unsure{+n|n|\prime n}}$ sont des bases de $E^{(2n)}$ et $E^{(-2n)}$ respectivement.

Donc
\[
L\ \unsure{\text{é}|\text{t}|{}^{\prime}}\ \text{Lefschetzien}\quad\Longrightarrow\quad L^n\ \text{\uline{une base de}}\ E^{2n}
\]
\inserted{(et si $\Lambda$\deleted{\unsure{\xi}} associé, $\bar\Lambda$ une base de $E^{-2n}$, \unsure{pour certains $L$})}

\margin{?} On \unsure{sait|saut} \unsure{penser|prouver|poser} une réciproque, dans \deleted{\unsure{le}} \inserted{certains} cas où \deleted{\unsure{$E$}} $E$ est projectif de type fini sur $k$. Ce \deleted{ne} marche certainement pas tel quel, par ex. quand $E=\mathcal L(M^*)$, $M^*$ gradué à degrés $0\leq i\leq4$ ---.

Mais ça ne marche pas non plus pour des anneaux associés \deleted{aux} éléments de degré \unsure{1|2} d'une alg. de Lefschetz, pour la dim 6

[ engendré par $L,L'$ \inserted{de degré\unsure{s|{}} 2}, avec les relations
\[
L^4=0,\ {L'}^{2}-LL'=0,\quad L^2L'-L^3=0,\ L^3L'=0\ ,
\]
donc base
\[
\begin{array}{c@{\ //\ }c@{\ //\ }c@{\ //\ }c}
1&L,L'&L^2,LL'&L^3\\ \hline
\text{degré }0&2&4&
\end{array}
\qquad\Big/\quad
\begin{aligned}
&L\ \text{est lefschetzien,}\\
&L'\ \text{ne l'est pas, bien que}\\
&{L'}^{3}=L^3\quad !\ .
\end{aligned}
\]
% archive sheet number: 21
\scanpage{23}\mspage{43}
\[
H^4(X,\underline{\mathrm Z}_l\handwritten{[}2\handwritten{]})\handwritten{\simeq}\underline{\mathrm Z}_l\ .
\]
\deleted{D'après \unsure{Il est}} Cette forme \inserted{(en négligeant la torsion du $H^2$) est non dégénérée, et} indu\deleted{\unsure{1}}it sur $\mathrm{NS}(X)\otimes_{\underline{\mathrm Z}}\underline{\mathrm Z}_l$ la forme bilinéaire déduite de la \uline{forme intersection} sur $\mathrm{NS}(X)$, provenant de la \deleted{\unsure{p}}multiplicité globale d'intersections de diviseurs sur $X$ \margin{\handwritten{$\times$}}\handwritten{[15]} \deleted{\unsure{;}}. D'après \margin{\handwritten{$\times$}}MATSUSAKA \handwritten{[27]}, cette dernière forme est non dégénérée. T\deleted{\unsure{a}}ensorisant al\unsure{o}rs la suite exacte (8.7) par $\underline{\mathrm Q}_l$, on tro\deleted{\unsure{n}}uve que $T_l(\mathrm{Br}^{\deleted{\unsure{\prime}}}(X))\otimes_{\underline{\mathrm Z}_l}\underline{\mathrm Q}_l$ \uline{s'identifie canoniquement à l'orthogonal de} $\mathrm{NS}(X)\otimes_{\underline{\mathrm Z}_l}\underline{\mathrm Q}_l$ \uline{dans} $H^2(X,\underline{\mathrm Q}_l\handwritten{[}1\handwritten{]})$, \uline{et que le cup-produit définit sur cet esp\deleted{\unsure{e}}ace une forme quadratique non dégénérée}. L'interprétant comme une forme quadratique sur $T_l(\mathrm{Br}^{\deleted{\prime}}(X))$ à valeurs dans $\underline{\mathrm Q}_l$, (et tenant compte que le cup-produit induit même une autodualité de $H^2/\text{torsion}$\unsure{)|),|;)} on trouve aisément que cette forme est à valeurs dans $\underline{\mathrm Z}_l$ si et seulement si le discriminant de la forme intersection sur $\mathrm{NS}(X)/\text{torsion}$ est une unité $l$-adique - condition qui est vérifiée donc pour presque tous les nombre premiers $l$ \unsure{-}, et que dans ce cas on trouve une autodualité de $T_l(\mathrm{Br}^{\deleted{\prime}}(X))$.

\handwritten{8.3.} Ce résultat ne concerne essentiellement que la partie divisible $\mathrm{Br}'(X)^{\unsure{o|0}}$ du groupe de Brauer cohomologique. On peut également donner une relation de dualité concernant le \deleted{\unsure{grou}} groupe fini $\mathrm{Br}'(X)/\mathrm{Br}'(X)^o$, qui grâce à la suite exacte de Kummer (BR II 3.1) s'explicite comme
\[
\mathrm{Br}'(X)/\mathrm{Br}'(X)^o\handwritten{\simeq}H^2(X,\handwritten{\mu}_{l^\infty})/H^2(X,\handwritten{\mu}_{l^\infty})^o\ , \tag*{(8.8)}
\]
(compte tenu que $\mathrm{Pic}(X)\otimes\underline{\mathrm Q}_l/\underline{\mathrm Z}_l=\mathrm{NS}(X)\otimes\underline{\mathrm Q}_l/\underline{\mathrm Z}_l$ est divisible). Utilisant (8.3) pour $i=3$, on trouve donc
\[
\mathrm{Br}'(X)/\mathrm{Br}'(X)^o\handwritten{\simeq}\mathrm{Tors}\ H^3(X,\underline{\mathrm Z}_l\handwritten{[}1\handwritten{]})\ , \tag*{(8.9)}
\]
\handwritten{(}valable indépendamm\deleted{\unsure{a}}ent des hypothèses particulières \deleted{\unsure{faites}}\inserted{posées} au début de l'alinéa précédent pour $X$. Lorsque celles-ci sont vérifiées, alors les relations de dualité,bien connues dans le cas classique, et qui résultent aisément \inserted{\handwritten{par passage à la limite}} des résultats de dualité \margin{\handwritten{$\times$}}\handwritten{[}SGAA \handwritten{$\underline{\overline{\mathrm{XIX}}}$\quad]} pour des
% archive sheet number: 22
\scanpage{25}
\diagram{\(\begin{array}{ccccc}
H^0&H^1&H^2&H^3&H^4\\
k&V&k\xi+W&\check V&k\xi^2
\end{array}\)}
\margin{Algèbres de Hodge d'ordre 2}

$V,W$
\[
\begin{aligned}
{\textstyle\bigwedge}^{2}V&\xrightarrow{\alpha}k\xi+W&
\alpha(x\wedge y)&=\varphi(x,y)\xi+\psi(x,y)&
\varphi(x,y)&\quad\text{forme alternée}\\
&&&&\psi(x,y)&\quad\text{$W$-\unsure{v|p}l. bil. alt. dans $W$}\\
\operatorname{Sym}^{2}(W)&\xrightarrow{\beta}k\xi^2&
\beta(xy)&=Q(x,y)\xi^2&
Q&\quad\text{forme bil. sym. sur $W$}
\end{aligned}
\]
\deleted{\unsure{Permet de recons}} La connaissance de $V,W,\varphi,\psi,\unsure{\bar Q|Q}$ permet de reconstituer $H^*$ à isom canonique près.
\[
\begin{aligned}
H^0\times H^i&\mathrel{\unsure{-|\to}}H^i:\quad\text{clair par}\quad 1.\unsure{x|\alpha}=x.1=x\qquad\text{tt }x.\\
H^1\times H^i&\longrightarrow H^{i+1}\quad\left|
\begin{aligned}
\deleted{\unsure{\text{d'où}|\text{dont}}}\ i=1\quad &xy=\varphi(x,y)\xi+\psi(x,y) &x,y\in V\\
i=2\quad &\left|\begin{aligned}
x\xi&=\xi x=\widetilde\varphi(x): &&\deleted{\unsure{(z)}}y\to\varphi(y,x)\\
xw&=wx=\rho(w,x):&&y\longrightarrow Q(\psi(y,x),w)
\end{aligned}\right. &\left|\begin{aligned}x&\in V\\y&\in V\\w&\in W\end{aligned}\right.\\
i=3\quad &xy'=-y'x=\deleted{\unsure{+}}\langle x,y'\rangle\xi^2 &x\in V,\ y'\in\check V\\
\deleted{i=4}\quad &\deleted{\unsure{x\xi^2=\xi}}
\end{aligned}\right.
\end{aligned}
\]
\[
H^2\times H^i\longrightarrow H^{i+2}\ \left|\quad j=2:\quad
\left|\begin{aligned}
\xi\cdot\xi&=\xi^2\\
\xi w=w\xi&=0\quad(w\in W)\\
ww'&=Q(w,w')\xi^2
\end{aligned}\right.\right.
\]
\margin{$(*)$ On pose $\widetilde\varphi:V\mathrel{\unsure{-|\to}}\check V$, avec $\langle x,\widetilde\varphi(y)\rangle=\varphi(x,y)$}

Quels que soient les choix de $\varphi,\psi$, on trouve une loi de composition anticommutative, admettant $1$ comme unité.

\uline{Condition d'associativité} s'écrit $(xy)(tz)=(tx)(yz)$, sous la forme
\[
\deleted{\unsure{\varphi}}\quad
\boxed{\varphi(x,y)\varphi(t,z)+Q(\psi(x,y),\psi(t,z))=\varphi(t,x)\varphi(y,z)+Q(\psi(t,x),\psi(y,z))}
\]
% slanted
\margin{exprime aussi que la forme $(xy|zt)$ en $xyzt$ est alternée en $\unsure{y,z|x,y,z|y^2}$ (donc en les \unsure{quatre lettres} !)}
\uline{Condition de Lefschetz}
\[
\begin{aligned}
\xi^{\unsure{2|?}}:H^0&\simeq H^4 &&\text{tjrs satisfait}\\
\xi:H^1&\longrightarrow H^3 &&\text{sss}\quad\boxed{\varphi\ \text{non dégénérée}}
\end{aligned}
\]
\uline{Condition de Poincaré}
\[
\begin{aligned}
H^4&\simeq k &&\text{tjrs satisfait}\\
H^0\times H^4&\longrightarrow H^4\ \text{et}\ H^1\deleted{\unsure{-}}\times H^3\mathrel{\unsure{-|\to}}H^4 &&\text{non dég : \uline{\unsure{itou|itrs|idem}}}\\
H^{\unsure{2|?}}\times H^2&\mathrel{-}H^{\unsure{4|2}} &&\text{non dég sss}\quad\boxed{Q\ \text{non dégénérée}}
\end{aligned}
\]
% archive sheet number: 24
\scanpage{26}
\deleted{\uline{Corrections à apporter à S.G.A.D. XV,XVI,XVII}.

\uline{Exp. XV}\quad Page 4. Après l\handwritten{'}énoncé de la proposition 1.2 , ajouter :

Voir (XVII Append. III, 1 ).

Page 6 ligne 13 : \textquotedblleft\ a\handwritten{$\vert$}été vu dans \deleted{\unsure{Exp.}} VII 8.4.\textquotedblright

Page 43. Ligne 11 : ajouter une flèche verticale dans le diagramme:
\diagram{\begin{tikzcd}[ampersand replacement=\&,row sep=small]
M(S'')\handwritten{\arrow[d]}\\ M(\overset{\handwritten{\scriptscriptstyle\wedge}}{S}'')
\end{tikzcd}}

Page 46. 5 lignes avant la fin :

Il résulte alors de \deleted{\unsure{Exp}}. X\handwritten{\unsure{,}}4.8\deleted{\unsure{K}}. a)

Page 60 : ligne 6. Remplacer (confer\hspace{2em} ) par (voir XVII \deleted{\unsure{,}}Append. III, 2 )

Page 61: 4 lignes avant la fin: ,$x$ étant une spécialisation de $x'$, 

Page 77: ligne 2: supprimer : de Exp. $\mathrm{VI}_{\mathrm B}$ 5.2. et

Ligne 10: après connexe, ajouter (XII 6.6. c) et XIV 4.8 et 4.8. bis ).

page 78: remplacer Remarque 6.3 par remarque 6.2.1.

page 81: 5 lignes avant la fin. Modifier le texte comme suit:

Mais alors $H_{\handwritten{\eta}}$ contient un sous-groupe de Cartan $C_{\handwritten{\eta}}$ de $G_{\handwritten{\eta}}$. On vient de voir que , quitte à restreindre $S$, $C_{\handwritten{\eta}}$ se prolonge en un sous-groupe de Cartan $C$ de $G$, que l'on peut supposer contenu dans $H$. Par hypothèse...

Page 87: 9 lignes avant la fin. Ajouter (Voir XVII\deleted{\unsure{,}} Append. III, 3 )

Page 86: 3 lignes avant la fin : auquel cas .

Page 87: 7 lignes avant la fin : EGA IV 8 et 6.3.

\margin{\handwritten{$\checkmark$}}Page 115: 10 lignes avant la fin: remplacer $C'$ par $T'$.

Page 123. Ligne 3 . Inscrire dans la parenthèse: (\deleted{\unsure{Exp.}}\inserted{\handwritten{voir}} XVII append. III\deleted{\unsure{,}} 2 )

Ligne 7. Remplacer $\dim H_t$ par $\dim L_t$.}
% archive sheet number: 25
\scanpage{27}
\uline{Cas particulier} Celui où $(xy)(zt)=0$ \inserted{[i.e. $(H^1)\deleted{\unsure{H|\otimes}}^{3}=0$ dans $H^4$ !]} quels que soient $x,y,z,t$, i.e. où l'image de $\alpha:\bigwedge^2 V\to k\xi+W$ est \deleted{\unsure{dans}} \deleted{\unsure{le|un}} ss-espace isotrope de $k(\xi)+W$. \deleted{Par exemple, si $W$ \unsure{a}, dans les alg. des\unsure{,} cela se réalise néc. ainsi : Si $W$ est \inserted{de dim} impaire, alors on peut \unsure{écrire}} $Q$ non dég. étant donnée, et supposant $W\ne0$, il existe (si $k$ alg. clos disons) une droite isotrope \inserted{$D$} dans $k\xi+W$ qui n'est pas dans $W$, et on prendra pour $\alpha$ \deleted{\unsure{une appl}} définie par une \deleted{forme}\inserted{\unsure{$\cdot W^L$|appl.}} bilinéaire \inserted{alt.} non dégénérée \deleted{sur} de $V\times V$ dans $D$ : alors $\varphi$ sera aussi non isotrope\ldots.

\begin{center}\rule{3em}{.4pt}\end{center}

Plus généralement, considérons la forme 4-linéaire alternée $\chi$ définie par
\[
xyzt=\chi(x,y,z,t)\xi^2\qquad (x,y,z,t\in V)
\]
Alors la forme \inserted{bilinéaire} induite par la \inserted{forme} $Q\oplus1$ (sur $\unsure{k|V}\xi+V$) sur $\bigwedge^{\unsure{2|?}}V$ par $\alpha$, qui définit une \deleted{forme} forme linéaire sur
\diagram{\begin{tikzcd}[ampersand replacement=\&,column sep=large]
\operatorname{Sym}^{2}(\bigwedge^2 V)\arrow[rr,bend right=18,"\text{\unsure{appl.|canon.}}"']\&\text{, se factorise en}\&\bigwedge^4 V
\end{tikzcd}}
moyennant $\chi$, et donc les données sont les suivantes

a) $V$ vect. (de dim finie) sur $k$

\deleted{b) $\varphi:\bigwedge^2 V\to k$}

b) $\chi:\bigwedge^4 V\to k$, définissant une forme bilinéaire symétrique $\chi'$ sur $\bigwedge^2 V$

c) un vectoriel $\widetilde W$, avec une forme \deleted{quadratique}\inserted{bil. sym.} $\widetilde Q$, et une application linéaire
\[
\bigwedge^2 V\xrightarrow{\alpha}\widetilde W
\]
compatible avec $\chi',Q$ [on dira que $\widetilde W,\widetilde Q$ est une \uline{extension} de la forme $\chi'$]

d) un élément $\xi$ de $\widetilde W$, tel que $Q(\xi,\xi)=1$.

\uline{Conditions de non \unsure{singularité|dégénérescence}}
\[
\left|\begin{aligned}
&Q\ \text{non dégénérée}\\
&\varphi(\unsure{x|s},y)=Q(\alpha(x\wedge y),\xi)\quad\text{non dégénérée}
\end{aligned}\right.
\]
% archive sheet number: 26
\scanpage{28}
[14]\quad \rule{4.7em}{.4pt}, \uline{Sur} \uline{quelques} \uline{points} \uline{d'algèbre} \uline{homologique}, Tohoku Math. J., 9 (1957), 119-221.

[15]\quad Hochschild, G., \uline{Local} \uline{class} \uline{field} \uline{theory}, Ann. of Math., 51 (1950), 331-347.

[16]\quad Hochschild, G. and Nakayama, T., \uline{Cohomology} \uline{in} \uline{class} \uline{field} \uline{theory}, Ann. of Math., 55 (1952), 348-366.

[17]\quad Hochschild, G. and Serre, J.-P., \uline{Cohomology} \uline{of} \uline{group} \uline{extensions}, T.A.M.S., 74 (1953), 110-134.

[18]\quad Lang, S., Abelian Varieties, Interscience Tracts No. 7, 1959, New York.

[19]\quad Nakayama, T., \uline{Cohomology} \uline{of} \uline{class} \uline{field} \uline{theory} \uline{and} \uline{tensor} \uline{product} \uline{modules}: I, Ann. of Math., 65 (1957), 255-267.

[20]\quad Ogg, A. P., \uline{Cohomology} \uline{of} \uline{abelian} \uline{varieties} \uline{over} \uline{function} \uline{fields}, Ann. of Math., 76 (1962), 185-212.

[21]\quad Rosenberg, A. and Zelinsky, D., \uline{Amitsur's} \uline{complex} \uline{for} \uline{inseparable} \uline{fields}, Osaka Math. J., 14 (1962), 219-240.

[22]\quad Serre, J.-P., Corps Locaux, Actualités scientifiques et industrielles No. 1296, 1962, Paris.

[23]\quad \rule{4.7em}{.4pt}, \textquotedblleft Dimension cohomologique\textquotedblright, Notes\unsure{.|{}} at Collège de France, 1962.

[24]\quad \rule{4.7em}{.4pt}, Groupes Algébriques et Corps de Classes, Actualités scientifiques et industrielles No. 1264, 1959, Paris.

[25]\quad \rule{4.7em}{.4pt}, Sur la topologie des variétés algébriques en caractéristique p, Symposium Int. de Top. Alg., 1958, Mexico City.

[26]\quad Shatz, S., \uline{Cohomology} \uline{of} \uline{artinian} \uline{group} \uline{schemes} \uline{over} \uline{local} \uline{fields}, Ann. of Math., to appear.

[27]\quad Tate, J., \uline{Cohomology} \uline{of} \uline{abelian} \uline{varieties} \uline{over} p-\uline{adic} \uline{fields}, (Cf. Notes by S. Lang, Princeton University, May, 1959.)
\mspage{45}
% archive sheet number: 27
\scanpage{29}
\deleted{\uline{Exemple} Prenons \unsure{$\varphi$ non dég|pour \illegible{1} dég}}
\[
\left|\begin{aligned}
\varphi:\bigwedge^2 V&\to k\qquad\text{non dégénérée}\\
&\inserted{\text{i.e. }\varphi\in\bigwedge^2\check V}\\
&\deleted{\unsure{\text{d'où }}\check\varphi\in\unsure{\bigwedge}}\\
\chi&=\varphi^2\in\bigwedge^4\check V\qquad \chi:\bigwedge^4 V\to k
\end{aligned}\right.
\]
\deleted{Alors $\chi$ comme forme sur $\bigwedge^2 V$ est \unsure{non}}

\uline{Ex} Prenons $\chi'$ non dégénérée, \deleted{\unsure{$\alpha$ un isom.}} $\alpha=\mathrm{id}$ i.e. $\widetilde W=\bigwedge^2 V$.

Alors un choix convenable de $\xi\in\unsure{\bar W|\widetilde W}=\bigwedge^2 V$ donne n'importe quelle $\varphi$ voulue, \unsure{au vecteur scalaire|au facteur scalaire|à un facteur scalaire}, donc n'importe quelle $\check\varphi$ \inserted{$\in\bigwedge^2 V$} non dégénérée, au \unsure{vecteur scalaire|facteur scalaire}. Considérons alors $\check\varphi^{\unsure{2|1|{}}}$ comme élément de \deleted{$\bigwedge^2 V\otimes$} l'algèbre $\bigwedge^{\unsure{*|\bullet|\vee}}V\otimes\bigwedge^\bullet V$ [on écrira $\bigwedge^2 V\subset V\otimes V$] et $\check\varphi^{2}$ comme élément de $\bigwedge^2 V\otimes\bigwedge^2 V\subset\bigwedge^*V\otimes\bigwedge^*V$ ; le carré de $\check\varphi$ dans $\bigwedge^*V$ \inserted{(qui est $\in\bigwedge^4 V$)} est l'image de $\check\varphi^{\unsure{(2)|(n)}}$ précédent par $\bigwedge^2 V\otimes\bigwedge^2 V\longrightarrow\bigwedge^4 V$. Utilisant la forme $Q\unsure{:|=}Q_{\chi'}$ définie par $\chi'$, on voit que $\check\varphi^{(2)}$ définit un endom de $\bigwedge^2 V$, \uline{soit $u$}. Je dis que l'orbite sous $u$ des \deleted{\unsure{él}}éléments $\xi$ et $\deleted{\unsure{d}}\check\varphi$ de $\bigwedge^2 V$ peuvent être infinis, si $\dim V\geq4$ (ou 6?) ($\unsure{W|V}$ est de dim paire) \deleted{\illegible{1}} et \inserted{les espaces engendrés par les uns et \deleted{\unsure{pl}} les autres} de dim non bornée (pour dim $V$ croissante\ldots).
% archive sheet number: 28
\scanpage{30}
\[
H^r(A,G)=(0)\quad\text{\uline{for}}\quad r>2\ .
\]

Proof. The duality theorem of [26] shows that $H^r(k,G')=(0)$ for $r>2$ if $G'$ is any AGS over $k$. Prop. 14 yields
\[
H^r(A,G)=(0)\quad\text{for}\quad r>3\ .
\]
To treat the case $r=3$, Th. 1 shows that we may assume $A$ is the field $k$. By Th. 4 and the duality theorem used above,
\[
H^3(k,G)\cong H^3(k,G_s)\ .
\]
However, the right hand group vanishes by results of Tate [27], and the proof is complete.

One may obtain additional information concerning the cohomology groups with coefficients in $\underset{\handwritten{\scriptscriptstyle\sim\!\sim}}{\mathrm{G}}_m$ by combining Th. 6 with results of Berkson, and Rosenberg-Zelinsky [5, 21]. Let $k$ be a field, and let $K$ be a purely inseparable extension of $k$. Let $\handwritten{\otimes}^{r}K$ denote the $k$-algebra obtained by forming the tensor product of $K$ with itself (over $k$) $r$-times. The face operators
\[
\handwritten{\otimes}^{r}K\to\handwritten{\otimes}^{r+1}K
\]
induce group homomorphisms
\[
H^s(\handwritten{\otimes}^{r}K,\underset{\handwritten{\scriptscriptstyle\sim\!\sim}}{\mathrm G}_m)\to H^s(\handwritten{\otimes}^{r+1}K,\underset{\handwritten{\scriptscriptstyle\sim\!\sim}}{\mathrm G}_m)\ .
\]
By combining the face operators in the obvious fashion, we obtain a complex, $C(s)$ for each $s>0$.
\[
C(s):H^s(K,\underset{\handwritten{\scriptscriptstyle\sim\!\sim}}{\mathrm G}_m)\to H^s(K\handwritten{\otimes}K,\underset{\handwritten{\scriptscriptstyle\sim\!\sim}}{\mathrm G}_m)\to H^s(K\handwritten{\otimes}K\handwritten{\otimes}K,\underset{\handwritten{\scriptscriptstyle\sim\!\sim}}{\mathrm G}_m)\to\ldots\ .
\]

Proposition 15. \uline{For} \uline{each} $s>0$, \uline{the} \uline{complex} $C(s)$ \uline{is} \uline{acyclic}.
\mspage{42}
% archive sheet number: 29
